\section{Surface examples supporting section 1}

These examples verify the complete subdivision definition, including all
characteristic-map conditions. The ambient surfaces are Hausdorff. This is
supporting work for the paper's Figure~2, not a replacement for the general
theory. Lean names are in the \texttt{GuibasStolfi} namespace.

\subsection{Shared square geometry}

\begin{lemma}[Disk--square parametrization]\label{lem:disk-square}
  %\lean{GuibasStolfi.Square.parametrization_exists, GuibasStolfi.Square.image_circle}

  There is a homeomorphism $d:\mathbb C\to\mathbb R^2$ taking the open unit
  disk onto $(0,1)^2$, the closed unit disk onto $[0,1]^2$, and the unit circle
  onto the square boundary.
\end{lemma}
\begin{proof}

  Apply Mathlib's convex gauge-rescaling homeomorphism theorem to the compact
  convex square with nonempty interior, then invert the resulting map.
  The boundary equality follows by taking the difference of the two disk images.
\end{proof}

\begin{definition}[Square boundary path]\label{def:square-path}
  \uses{lem:disk-square, def:simple-path}
  %\lean{GuibasStolfi.Square.path, GuibasStolfi.Square.diskToSquare_on_elem}

  Pull the four corners and four open sides, in cyclic order, back through
  $d|_{S^1}$. They form an eight-piece simple path. The restriction of $d$ to
  each piece maps homeomorphically onto its corner or whole open side.
\end{definition}
The chosen $d$ is noncomputable and is not assigned an orientation. The torus
orientation proof uses separate quotient charts.

\subsection{A sphere with one equatorial loop}

\begin{definition}[Sphere model and cells]\label{def:loop-sphere}
  \uses{def:point-line-disk}
  %\lean{GuibasStolfi.LoopSphere.sphereHomeomorph, GuibasStolfi.LoopSphere.vertex,
  %      GuibasStolfi.LoopSphere.edge, GuibasStolfi.LoopSphere.north,
  %      GuibasStolfi.LoopSphere.south}

  Use the one-point compactification $\mathbb C\cup\{\infty\}$, identified
  with the unit sphere in $\mathbb R^3$. Choose $p\in S^1$. The cells are
  $\{p\}$, $S^1\setminus\{p\}$, the open unit disk, and
  $\{z:|z|>1\}\cup\{\infty\}$.
\end{definition}

\begin{proposition}[Sphere subdivision]\label{prop:loop-sphere-subdivision}
  \uses{def:loop-sphere, def:subdivision, def:simple-path}
  %\lean{GuibasStolfi.LoopSphere.subdivision, GuibasStolfi.LoopSphere.vertices_count,
  %      GuibasStolfi.LoopSphere.edges_count, GuibasStolfi.LoopSphere.faces_count}

  These cells form a subdivision with one vertex, one loop edge, and two faces.
\end{proposition}
\begin{proof}
  \uses{def:endpoints, lem:endpoints-card}

  Extend $z\mapsto1/\overline z$ by exchanging $0$ and $\infty$.
  This involutive homeomorphism fixes $S^1$ pointwise. Use the disk inclusion
  as the first face map and its composite with inversion as the second.
  The one-point, one-arc circle partition works for both faces. Local plane
  charts, with an inverted chart at infinity, prove the manifold condition.
  Partition, shapes, and S4 follow from these maps. The endpoint theorem makes
  the unique vertex the loop's unique endpoint.
\end{proof}

\begin{corollary}\label{cor:loop-sphere-connected}
  \uses{prop:loop-sphere-subdivision, thm:graph-connected-iff}
  %\lean{GuibasStolfi.LoopSphere.graph_connected, GuibasStolfi.LoopSphere.connected}

  The sphere example has connected graph and connected ambient space.
\end{corollary}
\begin{proof}

  Its graph has one vertex. Apply the general connectivity theorem and
  nonemptiness of the sphere.
\end{proof}
This topological construction is separate from the finite algebra
\texttt{EdgeAlgebra.sphere}; their geometric relationship is not formalized.

\subsection{The one-face torus}

\begin{definition}[Torus model and cells]\label{def:torus-cells}
  %\lean{GuibasStolfi.Torus.standardTorus, GuibasStolfi.Torus.vertex,
  %      GuibasStolfi.Torus.horizontal, GuibasStolfi.Torus.vertical, GuibasStolfi.Torus.face}

  Let $T=(\mathbb R/\mathbb Z)^2$, homeomorphic to Mathlib's
  two-dimensional unit torus. Its vertex is $(0,0)$, its two open edges
  are the punctured coordinate circles, and its face is their complement.
  Write $q_T(x,y)=([x],[y])$.
\end{definition}

\begin{proposition}[Torus subdivision]\label{prop:torus-subdivision}
  \uses{def:torus-cells, def:square-path, def:subdivision}
  %\lean{GuibasStolfi.Torus.subdivision, GuibasStolfi.Torus.vertices_count,
  %      GuibasStolfi.Torus.edges_count, GuibasStolfi.Torus.faces_count,
  %      GuibasStolfi.Torus.edges_are_loops}

  The torus cells form a subdivision with one vertex, two loop edges, and one face.
\end{proposition}
\begin{proof}
  \uses{lem:endpoints-card}

  Reduction modulo one gives a homeomorphism from $(0,1)$ onto the punctured
  circle. Its product maps the open square homeomorphically onto the face.
  Translating that open disk gives a disk neighborhood at every torus point.
  The characteristic map $q_T\circ d$ maps the closed disk onto the whole
  torus, which is the face closure. Opposite open sides map homeomorphically
  onto the same whole edge, and all corners map to the vertex. Use the
  square boundary path for S4. Partition uniqueness and the general endpoint
  theorem give the loop statements and counts.
\end{proof}

\subsection{The Klein bottle quotient}

\begin{definition}[Klein bottle model]\label{def:klein-quotient}
  %\lean{GuibasStolfi.KleinBottle.flip, GuibasStolfi.KleinBottle.relation,
  %      GuibasStolfi.KleinBottle.projection_eq_iff, GuibasStolfi.KleinBottle.flip_ne}

  Put $C=(\mathbb R/2\mathbb Z)\times(\mathbb R/\mathbb Z)$ and
  $\tau(u,v)=(u+[1],-v)$. This is a fixed-point-free involutive homeomorphism.
  Define $K=C/{\sim}$, where $a\sim b$ iff $b=a$ or $b=\tau(a)$.
  For the quotient $\pi:C\to K$, equality $\pi(a)=\pi(b)$ is exactly this relation.
\end{definition}

\begin{lemma}[Square gluing and local charts]\label{lem:klein-charts}
  \uses{def:klein-quotient}
  %\lean{GuibasStolfi.KleinBottle.horizontal_gluing, GuibasStolfi.KleinBottle.vertical_gluing,
  %      GuibasStolfi.KleinBottle.rectangle_embedding, GuibasStolfi.KleinBottle.canonical_surjective,
  %      GuibasStolfi.KleinBottle.isTwoManifold}

  The map $q_K(x,y)=\pi([x]_2,[y]_1)$ satisfies
  \[
    q_K(x,1)=q_K(x,0),\qquad q_K(1,y)=q_K(0,-y).
  \]
  It is an open embedding on every open rectangle
  $(a,a+1)\times(b,b+1)$. The half-open square $[0,1)^2$ represents all points.
  The quotient is Hausdorff and has disk neighborhoods at every point.
\end{lemma}
\begin{proof}

  Saturation of an open set under the relation is its union with its
  $\tau$-image, so the projection is open. The relation is closed, giving the
  Hausdorff quotient. Equality of plane images forces equality of first
  coordinates modulo one; in a half-open interval of length one these real
  coordinates agree. The $\tau$ alternative is then impossible since
  $[1]_2\ne0$, and the second coordinates agree as well. This proves rectangle
  injectivity; openness and continuity supply the embeddings. Reducing a point
  of $C$ into $[0,2)\times[0,1)$ and applying $\tau$ when necessary gives
  the half-open unit-square representatives. Center an embedded rectangle
  at a representative to obtain its disk neighborhood.
\end{proof}

\begin{proposition}[Klein bottle subdivision]\label{prop:klein-subdivision}
  \uses{lem:klein-charts, def:square-path, def:subdivision}
  %\lean{GuibasStolfi.KleinBottle.subdivision, GuibasStolfi.KleinBottle.vertices_count,
  %      GuibasStolfi.KleinBottle.edges_count, GuibasStolfi.KleinBottle.faces_count,
  %      GuibasStolfi.KleinBottle.edges_are_loops}

  The $q_K$-images of $\{(0,0)\}$, $(0,1)\times\{0\}$,
  $\{0\}\times(0,1)$, and $(0,1)^2$ form a subdivision with one vertex,
  two loop edges, and one disk face.
\end{proposition}
\begin{proof}
  \uses{lem:endpoints-card}

  Half-open representatives give the partition. The local charts give
  the edge and face shapes. The characteristic map is $q_K\circ d$.
  The open square maps homeomorphically to the dense face, while its closed
  square maps onto $K$. Top and bottom sides agree; the right side maps to
  the whole left edge by $y\mapsto1-y$, using period one. All corners agree.
  These facts establish S4. The general endpoint theorem gives both loops.
\end{proof}

\begin{proposition}[Isomorphic underlying graphs]\label{prop:torus-klein-graph}
  \uses{prop:torus-subdivision, prop:klein-subdivision, def:graph-iso}
  %\lean{GuibasStolfi.KleinBottle.torusGraphIso}

  The torus and Klein bottle examples have isomorphic graphs.
\end{proposition}
\begin{proof}

  Match their unique vertices, horizontal edges, and vertical edges. Both
  edges in each graph are loops at that vertex, so all link relations agree.
\end{proof}

\subsection{Clara L\"oh's finite triangulation}

\begin{definition}[Credited combinatorial reference]\label{def:klein-triangulation}
  %\lean{GuibasStolfi.KleinBottleTriangulation.triangles,
  %      GuibasStolfi.KleinBottleTriangulation.simplices}

  Take the 18 Klein bottle triangles from Clara L\"oh's Lean 3 project
  \emph{Exploring Formalisation}, and their finite downward closure on
  $\{0,\ldots,8\}$, including the empty face.
\end{definition}
The exact source file and pinned revision are:
\begin{quote}\small
\texttt{src/gen\_simplicial\_complex.lean}\\
\texttt{a9ebece93862cc166fcf62d5c6ae0964ab571deb}
\end{quote}
The adaptation, original copyright, and Apache 2.0 license are recorded in
\texttt{ACKNOWLEDGMENTS.md} and \texttt{third\_party/exploring-formalisation/}.
The quotient construction above is separate; no realization homeomorphism
has been proved.

\begin{lemma}[Finite counts and incidence]\label{lem:klein-triangulation-counts}
  \uses{def:klein-triangulation}
  %\lean{GuibasStolfi.KleinBottleTriangulation.subset_closed,
  %      GuibasStolfi.KleinBottleTriangulation.vertices_count,
  %      GuibasStolfi.KleinBottleTriangulation.edges_count,
  %      GuibasStolfi.KleinBottleTriangulation.triangles_count,
  %      GuibasStolfi.KleinBottleTriangulation.every_edge_has_two_triangles}

  The simplex set is downward closed, with 9 vertices, 27 edges, and 18
  triangles. Each edge belongs to exactly two listed triangles.
\end{lemma}
\begin{proof}

  Downward closure follows from the powerset construction. The finite counts
  and edge-incidence statements are verified by kernel-checked decisions.
\end{proof}
